\section{Lesson 1: What Is Being Said?}

Before asking whether a metaphysical argument succeeds, you have to know what sentence is on the table. Serious philosophical disputes go off the rails right here more often than anywhere else: a reader skims a dense paragraph of Aquinas, turns it into a brisk summary, and then praises or refutes a proposition Aquinas never asserted.

Reading begins with two questions: \textbf{Claim} (what exactly is being asserted?) and \textbf{Meaning} (what do the important terms mean here?). Under \textbf{Claim}, every analysis gives two things: the author's own words with every qualification intact, and a \textbf{Plain-language reading} that removes unnecessary school diction without altering what is claimed. Along the way, we begin practicing a distinction that governs the rest of our work: fixing what a term means is not discovering a fact about the world, yet clear definitions are indispensable before any discovery can be tested.

\subsection{Separating the Proposition from the Sentence}

A grammatical sentence on the page is a vehicle; the \textbf{proposition} is the exact assertion it delivers. Classical prose compresses heavily, so extracting the proposition requires checking four features every time:

\begin{itemize}
\item \textbf{Subject:} What is the sentence speaking about---an individual thing, a shared nature, a causal series, or a proposition?
\item \textbf{Predicate:} What is affirmed or denied of that subject?
\item \textbf{Quantity:} Does the claim cover \emph{all} instances, \emph{some}, \emph{none}, or one individual? English translations of scholastic Latin often hide universal claims inside bare singulars: ``matter is in potentiality'' speaks of all matter as such.
\item \textbf{Qualifications:} What restrictions bound the assertion? Classical writers constantly narrow claims using three pairs: \emph{simpliciter} (unqualifiedly, absolutely speaking) against \emph{secundum quid} (in a certain respect); \emph{per se} (of itself, essentially) against \emph{per accidens} (incidentally); and \emph{inquantum} (``insofar as,'' isolating the exact aspect under which a predicate holds).
\end{itemize}

Dropping a qualification easily produces a fluent, plausible paraphrase that asserts a different claim. Watch this happen on a sentence from Aquinas, \emph{Summa Theologiae} I, q.\,3, a.\,1, \emph{corpus} (1920 English Dominican translation):

\begin{quote}
For although in any single thing that passes from potentiality to actuality, the potentiality is prior in time to the actuality; nevertheless, absolutely speaking, actuality is prior to potentiality; for whatever is in potentiality can be reduced into actuality only by some being in actuality.
\end{quote}

Drop the qualifiers---``in any single thing that passes from potentiality to actuality,'' ``in time,'' and ``absolutely speaking''---and you get a brisk paraphrase many readers produce without noticing: ``Actuality always comes before potentiality.'' That paraphrase is smooth, grammatical, and completely different from what Aquinas wrote---indeed, Aquinas rejects it in the very same sentence. An acorn precedes the oak tree; cold wood precedes hot wood. Within a single changing subject, potentiality comes first \emph{in time}. What Aquinas asserts is that \emph{absolutely speaking} (\emph{simpliciter})---in the causal order as a whole, since nothing potential activates itself without something already actual---actuality is prior. Strip away ``in time'' and ``absolutely speaking,'' and you turn a carefully bounded thesis into an absurd denial of ordinary growth.

The qualifier \emph{inquantum} (``insofar as'') does equally heavy work. In \emph{Summa Theologiae} I, q.\,46, a.\,2, \emph{ad 7}, Aquinas writes that a man begets \emph{insofar as} he is a man, and not \emph{insofar as} he is another man's son. Drop ``insofar as,'' and the sentence appears to deny that this man needed a father to exist. Kept intact, \emph{inquantum} separates what brought an agent into being from the nature by which he acts now. Whenever you meet ``insofar as,'' ``absolutely speaking,'' ``in a certain respect,'' ``\emph{per se},'' or ``\emph{per accidens},'' mark it with a pencil; if your restatement drops it, you have changed the claim.

\subsection{Technical and Ordinary Words, Decided Per Occurrence}

No reader can memorize a scholastic glossary once and read on autopilot. Six recurring words---\textbf{act}, \textbf{potency}, \textbf{form}, \textbf{subject}, \textbf{accident}, and \textbf{being}---can wear a technical sense in one sentence and an ordinary or shifted sense three lines later. You decide their sense \emph{per occurrence}:

\begin{itemize}
\item \textbf{Act} (\emph{actus}): Not a deed or theatrical performance, but realized existence or completed perfection. Yet elsewhere Aquinas uses ``act'' for an operation flowing from a power, as when he writes that ``the act of the eye is to see'' (\emph{ST} I, q.\,3, a.\,1, \emph{ad 3}).
\item \textbf{Potency} (\emph{potentia}): Not vigor or chemical strength. Scholastic usage separates \emph{active potency} (an agent's power to act, as fire heats) from \emph{passive potency} (a subject's real capacity to receive a perfection it currently lacks, as cold wood can become hot).
\item \textbf{Form} (\emph{forma}): Not mere geometric outline, but the determining principle making a thing the kind of thing it is (\emph{substantial form}, as the soul in a living body) or modifying an existing thing (\emph{accidental form}, as whiteness or heat).
\item \textbf{Subject} (\emph{subiectum}): In logic, what a proposition speaks about; in metaphysics, the underlying bearer in which a form or accident inheres (as bronze bears a statue's shape).
\item \textbf{Accident} (\emph{accidens}): Never a mishap. In metaphysics, a real feature that exists in another rather than on its own (height, color, knowledge); in logic, a feature outside a thing's definition.
\item \textbf{Being} (\emph{ens} or \emph{esse}): Its reference must be checked every single time. Does ``being'' here mean (1) the real act of existing (\emph{esse}), (2) an existing thing in the ten categories (a substance or real accident), or (3) the truth of an affirmative proposition?
\end{itemize}

Aquinas himself checks ``being'' per occurrence. In Chapter 1 of \emph{De ente et essentia} (Corpus Thomisticum Latin text), citing Aristotle's \emph{Metaphysics} V, he notes that being is said in two ways (\emph{ens per se dicitur dupliciter}): first, of what divides into the ten categories and posits something real; second, of what signifies the truth of propositions, under which even privations are called ``beings,'' as in ``blindness is in the eye.'' Blindness has no essence and adds no quality; it is the absence of sight. Miss which sense is active, and you mistake a true report of an absence for an existing thing.

Fixing \textbf{Meaning} obeys two rules. Define only terms that are technical, active in the inference, or ambiguous; an illustration needs no essay. Fix those senses \emph{before} evaluating the argument and keep them locked. If a word shifts sense between premises, mark both readings at once: an argument true on one reading and valid on another is equivocal, not sound. Never bend a definition mid-analysis to rescue an inference.

\subsection{Rules for a Plain-Language Reading, and Testing Distinctions}

The \textbf{Plain-language reading} under \textbf{Claim} proves whether you grasped what the author said, and it is also the easiest place to rewrite an argument accidentally. Five rules govern it:

\begin{enumerate}
\item \textbf{Simplify diction, never the claim.} Drop unnecessary school jargon, never a real distinction.
\item \textbf{Never strengthen or weaken.} Never turn ``some'' into ``all,'' ``is'' into ``must,'' ``not judged impossible'' into ``proved possible,'' or ``in this respect'' into ``unqualifiedly.''
\item \textbf{Never add or move a quantifier.} ``Every moving thing has some mover'' is not ``one mover moves every moving thing.''
\item \textbf{Never silently resolve an open ambiguity.} If a sentence legitimately admits two readings, mark the ambiguity explicitly instead of picking the tidier one.
\item \textbf{Keep what cannot be replaced without loss.} Where a term like \emph{per se} or \emph{act} carries a precise sense no single everyday word captures, keep it, explain it right in the sentence, and say why flattening it would alter the claim.
\end{enumerate}

Test every plain-language reading by one question: \textbf{would the author accept it as an accurate statement of his own position?} If Aristotle or Aquinas could fairly answer, ``You dropped my qualifier,'' or ``You made me claim more than I wrote,'' the restatement fails, however fluent it reads.

Fixing meaning also teaches you to judge a scholastic \textbf{distinction}. Trained readers do not dismiss every \emph{distinguo} (``I distinguish'') as verbal dodging, nor do they accept every distinction on credit. Apply the \textbf{independent-identification test}: a distinction genuinely answers an objection only if both senses can be identified \emph{independently of the objection}, using examples or criteria that hold outside the dispute.

Watch Aquinas pass that test in \emph{Summa Theologiae} I, q.\,3, a.\,4, \emph{ad 1} (1920 English Dominican translation). Objection 1 argued that if God's essence and existence are identical, divine being has ``nothing added to it''; yet being with nothing added to it is universal being predicated of stones and wood, making God identical with general being. Aquinas replies:

\begin{quote}
A thing that has nothing added to it can be of two kinds. Either its essence precludes any addition; thus, for example, it is of the essence of an irrational animal to be without reason. Or we may understand a thing to have nothing added to it, inasmuch as its essence does not require that anything should be added to it; thus the genus animal is without reason, because it is not of the essence of animal in general to have reason; but neither is it to lack reason. And so the divine being has nothing added to it in the first sense; whereas universal being has nothing added to it in the second sense.
\end{quote}

Both senses of ``without addition''---excluding addition versus leaving addition open---stand cleanly on ordinary logical cases (``irrational animal'' against ``animal in general'') that anyone recognizes before God is ever mentioned. The distinction exposes a real ambiguity in ordinary speech. Contrast an ad-hoc evasion: someone claims ``every metal sinks in water,'' sees sodium float, and replies, ``By `true metal' I mean only metals that sink.'' Because the two senses cannot be identified apart from dodging the counterexample, the distinction is merely a disguised restatement of the claim.

\subsection{Worked Passage 1: A Sentence Whose Truth Is Mostly Definitional}

Our first worked passage draws on Aristotle, \emph{Posterior Analytics} I.2, 71b17--25 (G.\ R.\ G.\ Mure's Oxford translation), where Aristotle stipulates: ``By demonstration I mean a syllogism productive of scientific knowledge,'' whose premises must be ``true, primary, immediate, better known than, prior to, and causes of the conclusion.''

Examine this sentence drawn from that passage: \emph{A valid syllogism from true premises may still fail to be a demonstration.} It sounds weighty and substantive. Let us take it through \textbf{Claim} and \textbf{Meaning}.

\subsubsection*{Claim}
From Aristotle's definition in \emph{Posterior Analytics} I.2 (71b17--25): a valid syllogism from true premises may still fail to be a demonstration.

\emph{Plain-language reading.} Even if an argument has a valid deductive form and every premise it uses is true, it does not count as a demonstration in Aristotle's sense unless its premises also satisfy his further requirements---being primary, immediate, prior in nature, better known in nature, and causes of the conclusion.

\subsubsection*{Meaning}
Three terms carry technical senses fixed in Aristotle's treatises, while two stay ordinary. \emph{Syllogism} (\emph{syllogismos}) is a deductive inference in which, certain things being laid down, something other than them follows of necessity (\emph{Prior Analytics} I.1, 24b18--20). \emph{Scientific knowledge} (\emph{episteme}) is knowing a necessary fact through its proper cause (\emph{Posterior Analytics} I.2, 71b8--16). \emph{Demonstration} (\emph{apodeixis}) is stipulated directly here: a syllogism productive of scientific knowledge, meeting all six conditions on its premises. \emph{Valid} and \emph{true} carry their ordinary logical senses, and no term shifts its meaning.

This sentence accomplishes something exact, yet makes no discovery about the furniture of the world. Once Aristotle stipulates that ``demonstration'' requires six conditions beyond validity and truth, it follows immediately from the definition alone that having only validity and truth leaves an argument short of demonstration. This is a standard, not a discovery about reality. Saying so is not belittling the sentence: fixing a crisp standard is essential before anyone can test whether a real argument meets it. No symbolic notation is used here, because no scope, identity, or formal ambiguity is at issue; symbols would merely dress up a definition.

\subsection{Worked Passage 2: Three Technical Terms and a Load-Bearing Qualification}

Now apply \textbf{Claim} and \textbf{Meaning} to a real sentence of St.\ Thomas Aquinas from \emph{Summa Theologiae} I, q.\,3, a.\,4, \emph{corpus} (``Whether Essence and Existence Are the Same in God?'', 1920 English Dominican translation). Here is the opening of the second argument in the \emph{corpus}:

\begin{quote}
Secondly, existence is that which makes every form or nature actual; for goodness and humanity are spoken of as actual, only because they are spoken of as existing. Therefore existence must be compared to essence, if the latter is a distinct reality, as actuality to potentiality.
\end{quote}

Focus on the second sentence: \emph{Therefore existence must be compared to essence, if the latter is a distinct reality, as actuality to potentiality.} It carries four technical terms---\emph{existence}, \emph{essence}, \emph{actuality}, and \emph{potentiality}---and one load-bearing qualification: \emph{if the latter is a distinct reality}. Test three candidate paraphrases against the original.

\subsubsection*{Candidate 1 (rejected)}
``So being real is to a thing's definition like action is to strength.'' This collapses all four technical terms into unrelated everyday cognates: ``being real'' blurs existence with factual truth, ``definition'' replaces a real essence with a verbal formula, and ``action'' and ``strength'' distort \emph{actuality} and \emph{potentiality}.

\subsubsection*{Candidate 2 (rejected)}
``In every being, existence stands to essence as actualizing perfection stands to receptive capacity.''

Candidate 2 is the trap serious readers fall into: grammatical, accurate in diction, and true of every \emph{created} thing. Yet as a reading of Aquinas's sentence, \textbf{it fails}: it drops the load-bearing word \textbf{if} (``if the latter is a distinct reality'') and puts ``In every being'' in its place. Two lines later in \emph{ST} I, q.\,3, a.\,4, \emph{corpus}, Aquinas invokes his prior claim that God has no potentiality and concludes that in God essence and existence are \emph{not} distinct. Aquinas's sentence is a conditional step in a proof by contradiction; dropping \textbf{if} turns it into an unqualified universal that contradicts his own conclusion.

\subsubsection*{Candidate 3 (accepted): Claim and Meaning}
\textbf{Claim.} As Aquinas writes (\emph{ST} I, q.\,3, a.\,4, \emph{corpus}): ``Therefore existence must be compared to essence, if the latter is a distinct reality, as actuality to potentiality.''

\emph{Plain-language reading.} Consequently, whenever a thing's essence is genuinely distinct from its act of existing, its act of existing stands to that essence as realized perfection stands to a receptive capacity awaiting realization.

Aquinas would accept this reading: it keeps the inferential marker, preserves the conditional restriction, and renders \emph{actuality} and \emph{potentiality} without implying that essence and existence differ in God.

\textbf{Meaning.} Four technical terms do the work: \textbf{existence} (\emph{esse}), the real act making a form or nature actual rather than mere propositional truth; \textbf{essence} (\emph{essentia}), the real nature making a thing the kind of thing it is; \textbf{actuality} (\emph{actus}), realized perfection; and \textbf{potentiality} (\emph{potentia}), passive receptive capacity open to an actuality not possessed of itself. The qualification \textbf{if the latter is a distinct reality} leaves open whether essence and existence differ in a given subject, asserting only that \emph{wherever} they differ, they stand as potency to act.

\subsection{Exercises}

\subsubsection*{Exercise 1.1}
Read this sentence from \emph{Summa Theologiae} I, q.\,3, a.\,3, \emph{corpus} (1920 English Dominican translation), with three technical expressions italicized:

\begin{quote}
in things composed of \emph{matter and form}, the nature or essence must differ from the \emph{suppositum}, because the essence or nature connotes only what is included in the definition of the \emph{species}; as, humanity connotes all that is included in the definition of man.
\end{quote}

Explain what each italicized expression means here, and write a plain-language reading preserving the opening restriction (``in things composed of matter and form'').

\subsubsection*{Exercise 1.2}
Read this sentence from \emph{Summa Theologiae} I, q.\,46, a.\,1, \emph{ad 2} (1920 English Dominican translation), with two qualifiers italicized:

\begin{quote}
Whatever has power always to be, \emph{from the fact of having that power}, cannot sometimes be and sometimes not be; but \emph{before it received that power}, it did not exist.
\end{quote}

Explain how the claim alters if the first italicized qualifier is dropped, and write a plain-language reading preserving both qualifiers.

\subsubsection*{Exercise 1.3}
Read this sentence from \emph{Summa Theologiae} I, q.\,44, a.\,1, \emph{ad 1} (1920 English Dominican translation):

\begin{quote}
Though the relation to its cause is not part of the definition of a thing caused, still it follows, as a consequence, on what belongs to its essence; because from the fact that a thing has being by participation, it follows that it is caused.
\end{quote}

A reader offers this paraphrase: ``Everything that exists by participation is caused by God, even though being caused is not part of its definition.'' The paraphrase is grammatical and states something Aquinas affirms in the same article. Explain why it fails as a plain-language reading of this sentence, name the exact words in the original doing the work that it drops or alters, and write a repaired plain-language reading.

\subsubsection*{Exercise 1.4}
Read Aquinas's reply in \emph{Summa Theologiae} I, q.\,3, a.\,4, \emph{ad 2} (1920 English Dominican translation), answering an objection that because we know \emph{whether} God exists without knowing \emph{what} God is, God's existence must differ from His essence:

\begin{quote}
``To be'' can mean either of two things. It may mean the act of essence, or it may mean the composition of a proposition effected by the mind in joining a predicate to a subject. Taking ``to be'' in the first sense, we cannot understand God's existence nor His essence; but only in the second sense. We know that this proposition which we form about God when we say ``God is,'' is true; and this we know from His effects (Q.\,2, A.\,2).
\end{quote}

State the two senses of ``to be'' Aquinas distinguishes, and test the distinction by the independent-identification rule: is it a real distinction or an ad-hoc evasion?

\subsubsection*{Exercise 1.5}
Read this unannotated paragraph from \emph{Summa Theologiae} I, q.\,3, a.\,1, \emph{corpus} (1920 English Dominican translation), containing Aquinas's third argument that God is not a body:

\begin{quote}
Thirdly, because God is the most noble of beings. Now it is impossible for a body to be the most noble of beings; for a body must be either animate or inanimate; and an animate body is manifestly nobler than any inanimate body. But an animate body is not animate precisely as body; otherwise all bodies would be animate. Therefore its animation depends upon some other thing, as our body depends for its animation on the soul. Hence that by which a body becomes animated must be nobler than the body. Therefore it is impossible that God should be a body.
\end{quote}

With nothing marked for you, give a complete \textbf{Claim} (with \textbf{Plain-language reading}) and \textbf{Meaning} analysis, naming the technical terms and the exact qualification doing the work.

\subsection{Solutions}

\subsubsection*{Solution to Exercise 1.1}
\emph{Matter and form} are the two constitutive principles of physical things---underlying material (this flesh and these bones) and the determining nature making them a thing of one kind (the soul). The \emph{suppositum} is the complete concrete individual that actually exists (Socrates himself, with his individual traits), as distinct from abstract humanity. The \emph{species} is the general kind captured by the definition (``human being''), excluding traits unique to one individual.

Plain-language reading: ``In physical substances composed of matter and form, a thing's shared nature differs from the concrete individual possessing it, because the shared nature includes only what belongs to the definition of its kind---just as humanity includes only what defines a human being as such.'' Dropping ``in things composed of matter and form'' ruins Aquinas's article, because two sentences later he concludes that in immaterial realities not composed of matter and form, individual and nature are identical.

\subsubsection*{Solution to Exercise 1.2}
Dropping ``from the fact of having that power'' turns the first clause into ``Whatever has power always to be cannot sometimes be and sometimes not be''---which concedes Objection 2's whole argument that an incorruptible creature could never have begun to exist! With the qualifier intact, Aquinas restricts the impossibility strictly to what follows \emph{once the thing possesses that enduring power}: having received an incorruptible nature, it cannot perish or flicker out. The second qualifier (``before it received that power'') then distinguishes an intrinsic incapacity to perish from having existed without a beginning: prior to receiving that nature by creation, it did not exist.

Plain-language reading: ``Once a thing possesses a nature capable of enduring forever, possessing that nature prevents it from alternating between existing and not existing; yet before it received that nature, it did not exist.'' Paraphrasing ``has power always to be'' as ``is eternal'' would smuggle in the very sense of beginningless existence Aquinas is excluding.

\subsubsection*{Solution to Exercise 1.3}
Although grammatical and true within Aquinas's theology, the paraphrase fails for two concrete reasons. First, it drops the load-bearing phrase \textbf{``still it follows, as a consequence, on what belongs to its essence.''} Those words supply the exact link answering Objection 1: between what lies \emph{inside} a definition and what is merely incidental stands a necessary consequence flowing from the essence (like the human capacity for laughter in Aquinas's next clause). Omitting those words erases the entire explanation of why something outside a creature's definition can still be inseparable from it. Second, the paraphrase inserts \textbf{``by God,''} which this sentence does not state; here Aquinas argues only that a participated being is \emph{caused}, leaving the identification of its cause to the surrounding article.

Repaired plain-language reading: ``Although standing in relation to its cause is not included inside a caused thing's definition, it still follows necessarily from what belongs to its nature: because a thing possesses existence only by receiving a share in it, being caused follows inseparably from what it is.''

\subsubsection*{Solution to Exercise 1.4}
Aquinas distinguishes two senses of ``to be'': (1) the real \emph{act of essence}---concrete existence in reality; and (2) the \emph{composition of a proposition} by which the mind joins predicate to subject, i.e., the truth of an affirmative statement such as ``God is.'' This passes the independent-identification test: both senses appear outside theology whenever we say ``blindness is in the eye'' (a true affirmative proposition, not a real substance) as against asking what a living tree's act of existing is. Because both senses are recognizable independently of Objection 2, distinguishing them clarifies rather than evades.

\subsubsection*{Solution to Exercise 1.5}
\textbf{Claim.} As Aquinas writes (\emph{ST} I, q.\,3, a.\,1, \emph{corpus}): God is the most noble of beings; no body can be the most noble, because an animate body surpasses an inanimate one yet is not animate precisely as body, so whatever animates it must surpass the body; hence God cannot be a body.

\emph{Plain-language reading.} God stands highest in excellence among all beings. No physical body can hold that place: every body is living or non-living, and a living body surpasses a non-living one, yet a living body is not alive merely by being a body (else every body would be alive). Its life therefore depends on another principle, as our body depends on the soul; and whatever gives a body life must surpass the body in excellence. Thus God cannot be a body.

\textbf{Meaning.} Four technical terms carry weight: \textbf{body} (\emph{corpus}), a three-dimensional material thing; \textbf{noble} (\emph{nobilis}), ontological perfection rather than social rank; \textbf{animate / inanimate}, living versus non-living (``animation'' means being alive); and \textbf{soul} (\emph{anima}), the formal life-principle of an organic body. The load-bearing qualification is \textbf{``precisely as body''} (Latin \emph{inquantum est corpus}); drop it, and ``a living body is not alive'' becomes a flat contradiction, whereas keeping it isolates the formal respect under which life does not belong to bodiliness as such. One further point bears watching later: ``nobler'' must keep one steady sense when comparing living to non-living bodies and when comparing the life-principle to the body itself.
